
@@M0@@
étant le morphisme structural @@M1@@.

Les @@M2@@ considérés seront ceux qui sont définis à l'aide des faisceaux localement libres de type fini. On notera @@M3@@, @@M4@@, @@M5@@, les classes de @@M6@@, @@M7@@, @@M8@@. On pose @@M9@@, @@M10@@, et on note @@M11@@, @@M12@@ les classes de @@M13@@, @@M14@@ dans @@M15@@. Comme @@M16@@ et @@M17@@ sont dans @@M18@@, on a @@M19@@, i.e. @@M20@@ mod. @@M21@@ i.e. @@M22@@

5.2.
Si @@M23@@ est un @@M24@@-Module localement libre de type fini, les @@M25@@ ne sont pas en général localement libres de type fini sur @@M26@@, si bien qu'on ne peut définir un homomorphisme @@M27@@ comme on le fait pour les @@M28@@ définis à l'aide des complexes parfaits. Le théorème 1.1 permet de tourner cette difficulté de la façon suivante : si @@M29@@, on peut, d'une seule manière, écrire :
@@M30@@
on pose alors :
@@M31@@

Il est clair que l'homomorphisme ainsi défini est additif, et vérifie la formule de projection, i.e. est @@M32@@-linéaire.

On vérifie les relations suivantes :
@@M33@@

Pour vérifier le dernier groupe de relations, on procède par récurrence sur @@M34@@. Supposons que pour @@M35@@, @@M36@@; la relation (1.9.1), multipliée par une puissance convenable de @@M37@@, donne :
@@M38@@
On peut lui appliquer @@M39@@ :
